% ============================================================================
%  ANEXOS TÉCNICOS DEL SUBDOCUMENTO 2
%  audIT Soluciones Tecnológicas SpA - Licitación TFEP-01/2026
%  Caso 10: Transportes Curimón S.A.
% ============================================================================

El presente documento complementario contiene los inventarios detallados, matrices de requerimientos, caracterización exhaustiva de flota y conductores, matrices normativas y fichas pormenorizadas de actores que respaldan analíticamente el \textbf{Subdocumento 2} (\textit{Comprensión del Problema y de la Necesidad}) de la oferta técnica.

Conforme a lo instruido en el Comunicado 10 (§1 y §5), las tablas de catalogación y los listados que superan cinco columnas o una página de extensión se presentan en este archivo de anexos independiente para salvaguardar la legibilidad del cuerpo principal. Este documento se estructura en cuatro anexos formales:
\begin{itemize}
  \item \textbf{Anexo 2.A:} Catálogo Exhaustivo de Requerimientos de Negocio y Operacionales Preliminares.
  \item \textbf{Anexo 2.B:} Inventario Detallado de Flota y Caracterización de Conductores.
  \item \textbf{Anexo 2.C:} Matriz Exhaustiva de Restricciones Operacionales, Legales y Exclusiones Contractuales.
  \item \textbf{Anexo 2.D:} Fichas Detalladas de Caracterización de los Trece (13) Actores del Ecosistema.
\end{itemize}

\vspace{0.5cm}

% ============================================================================
%  ANEXO 2.A
% ============================================================================
\section*{Anexo 2.A: Catálogo exhaustivo de requerimientos preliminares de negocio y operacionales}\label{sec:2-anexo-2a}
\addcontentsline{toc}{section}{Anexo 2.A: Catálogo exhaustivo de requerimientos preliminares de negocio y operacionales}

El catálogo compendia las necesidades preliminares de negocio levantadas desde las Bases Técnicas del Caso 10, las entrevistas en terreno con los actores operacionales y las exigencias normativas del transporte terrestre chileno. Cada requerimiento se codifica unívocamente, estableciendo su trazabilidad formal y su nivel de criticidad para la continuidad operacional en la \tabref{2-req-preliminares}.

\begin{tabla}[ancho=completo]{Catálogo de requerimientos de negocio y operacionales preliminares}{2-req-preliminares}{C{24mm} L{30mm} X L{32mm} C{18mm}}
\textbf{ID Req.} & \textbf{Dominio} & \textbf{Descripción Detallada del Requerimiento Preliminar} & \textbf{Fuente / Base} & \textbf{Criticidad} \\ \midrule
\textbf{REQ-NEG-01} & Asignación y Despacho & Validar síncronamente en pre-despacho ($\le 30\text{ s}$) que el conductor cuente con horas de jornada disponibles conforme al Art. 25 bis del Código del Trabajo, bloqueando la asignación si se superan las 5 h de manejo o no se acredita descanso previo. & Caso 10, Cap. 4.3; Entrevista R. Mansilla & Crítica \\
\textbf{REQ-NEG-02} & Asignación y Despacho & Cotejar automáticamente el estado de vencimiento de las ~6.000 vigencias vivas (licencias A5, revisiones técnicas, certificados de gases, permisos, seguros), impidiendo despachar vehículos o choferes con documentación caducada. & Caso 10, Cap. 4.4; Entrevista D. Aguayo & Crítica \\
\textbf{REQ-NEG-03} & Asignación y Despacho & Verificar la aptitud física del equipo asignado respecto al tipo de carga requerida (semirremolque refrigerado para perecibles, tolva para granel, o autorización D.S. N.° 298 para sustancias peligrosas). & Caso 10, Cap. 4.5; Entrevista R. Mansilla & Crítica \\
\textbf{REQ-NEG-04} & Sustancias Peligrosas & Comprobar de forma obligatoria que el conductor asignado a una de las 18 unidades SUSPEL cuente con el curso específico vigente del D.S. N.° 298 y que el vehículo porte Hoja de Datos de Seguridad y rotulación NCh 2190. & Caso 10, Cap. 4.5; Entrevista D. Aguayo & Crítica \\
\textbf{REQ-NEG-05} & Trazabilidad y Geocercas & Detectar automáticamente mediante geocercas poligonales la entrada, tiempo de permanencia y salida en los ~1.400 puntos de clientes, sin requerir intervención manual del conductor ni instalación de equipos en predios ajenos. & Caso 10, Cap. 4.7; Entrevista E. Valdebenito & Alta \\
\textbf{REQ-NEG-06} & Cobro de Sobreestadías & Evaluar registros cronológicos de llegada, espera y salida con datos de ubicación y tiempo para apoyar la revisión de cobros objetados. La integridad, precisión y admisibilidad de estos antecedentes deben validarse. & Caso 10, Cap. 4.7; Entrevista G. Ossandón & Alta \\
\textbf{REQ-NEG-07} & Retornos en Vacío & Identificar en tiempo real los tractocamiones que finalizarán su descarga para sugerir triangulaciones con cargas de retorno compatibles, reduciendo el 26\% de kilómetros recorridos en vacío (10,66 millones de km anuales). & Caso 10, Cap. 4.2; Entrevista R. Mansilla & Alta \\
\textbf{REQ-NEG-08} & Cadena de Frío & Monitorear en tiempo real la temperatura interna de las 44 ramplas refrigeradas (-30 °C a +30 °C), emitiendo alertas inmediatas a la Torre 24x7 ante desviaciones térmicas de $\pm 1{,}5\text{ }^\circ\text{C}$ o apertura no autorizada de puertas. & Caso 10, Cap. 2.1 y 4.8; Entrevista A. Lecaros & Crítica \\
\textbf{REQ-NEG-09} & Documentación Digital & Emitir Documentos Electrónicos de Transporte (DET para ~128.000 guías anuales) integrados con el ERP contable y el SII, habilitando la emisión offline pre-firmada en zonas de carga sin cobertura celular. & Caso 10, Cap. 4.6; Entrevista M. Riquelme & Alta \\
\textbf{REQ-NEG-10} & Confirmación de Entrega & Digitalizar el comprobante de entrega (\textit{Proof of Delivery} [POD]) mediante captura fotográfica y firma digital en pantalla en destino, abatiendo el 4,2\% de pérdidas o roturas de guías físicas en papel. & Caso 10, Cap. 4.7; Entrevista G. Ossandón & Media \\
\textbf{REQ-NEG-11} & Costeo por Ruta y Viaje & Reconstruir el costo marginal directo real de cada viaje (< 24 h post-cierre), integrando consumo de diésel por CAN bus, pasadas de peajes TAG y flete liquidado a terceros, erradicando el prorrateo ciego por ingreso. & Caso 10, Cap. 4.1 y 7.3; Entrevista G. Ossandón & Crítica \\
\textbf{REQ-NEG-12} & Renegociación Contratos & Proveer a la Gerencia de Finanzas la matriz de rentabilidad histórica desagregada por cliente y ruta para renegociar los 3 contratos deficitarios (31\% del ingreso, peor a -14\%) previo a sus vencimientos en 2027. & Caso 10, Cap. 2.3; Entrevista G. Ossandón & Crítica \\
\textbf{REQ-NEG-13} & Telemetría CAN bus & Capturar y procesar de forma pasiva y no intrusiva los parámetros de operación del motor (RPM, odómetro, temperatura de refrigerante, códigos DTC y consumo acumulado) en los 61 tractos con CAN bus de fábrica. & Caso 10, Cap. 4.10; Entrevista H. Trincado & Alta \\
\textbf{REQ-NEG-14} & Mantenimiento Preventivo & Generar órdenes automáticas de mantenimiento en base al kilometraje y horas de motor efectivamente acumulados por telemetría, sustituyendo la planificación visual manual cada 6 días en taller San Bernardo. & Caso 10, Cap. 4.10; Entrevista H. Trincado & Alta \\
\textbf{REQ-NEG-15} & Integración Talleres Ruta & Habilitar un canal web simplificado para que los talleres externos en ruta registren intervenciones de emergencia y repuestos instalados, actualizando la hoja de vida vehicular de forma inmediata. & Caso 10, Cap. 4.10; Entrevista H. Trincado & Media \\
\textbf{REQ-NEG-16} & Liquidación a Terceros & Automatizar el cálculo de pre-liquidaciones mensuales a los 148 transportistas terceros a partir de los viajes validados en sistema, reduciendo el ciclo de 9 días hábiles y la tasa de error del 11\%. & Caso 10, Cap. 4.11; Entrevista G. Ossandón & Alta \\
\textbf{REQ-NEG-17} & Privacidad de Terceros & Definir alcance, finalidad, acceso y retención de los datos de conductores y transportistas. Considerar la Ley N.° 19.628 vigente a la fecha de la oferta y preparar la adaptación a la Ley N.° 21.719, cuya entrada en vigor está prevista para el 1 de diciembre de 2026. La base jurídica requiere revisión. & Caso 10; Entrevista N. Sandoval & Crítica \\
\textbf{REQ-NEG-18} & Homologación Plataformas & Ingerir y unificar en una vista de mapa única las posiciones GPS provenientes de las tres plataformas dispares existentes (Wialon, Wisetrack, Webfleet) para los 192 camiones terceros que cuentan con rastreo previo. & Caso 10, Cap. 5; Entrevista P. Kast & Alta \\
\textbf{REQ-NEG-19} & Sensorización 34 Camiones & Equipamiento telemático estándar en los 34 camiones de terceros que carecen de GPS, incorporándolos a la vista operacional de la Torre de Control. & Caso 10, Cap. 2.1 y 5; Entrevista E. Valdebenito & Alta \\
\textbf{REQ-NEG-20} & Resiliencia Desconexión & Retención local mínima de 72 horas conforme a RT-03.10; cualquier capacidad superior propuesta para el hardware requiere dimensionamiento y pruebas de desconexión. & Caso 10, RT-03.10 & Crítica \\
\textbf{REQ-NEG-21} & Seguridad en Cabina & Evaluar controles de interfaz para reducir distracciones durante la conducción. Definir restricciones táctiles, alertas y condiciones de activación mediante revisión de la Ley N.° 21.377, su reglamentación y pruebas de seguridad. & Caso 10; Entrevista Y. Colipán & Crítica \\
\textbf{REQ-NEG-22} & Alerta Anticipada Fatiga & Calcular la alerta de descanso del Art. 25 bis considerando la distancia y tiempo estimado hacia el próximo punto seguro de detención (berma o servicentro), evitando que la alarma venza en zonas desérticas sin servicios. & Caso 10, Cap. 4.3; Entrevista Y. Colipán & Alta \\
\textbf{REQ-NEG-23} & Tacógrafo Digital & Evaluar la descarga y conservación de los archivos nativos del tacógrafo, con trazabilidad de acceso y controles de integridad. El formato exigible, periodicidad y valor probatorio se deben confirmar con la normativa aplicable y la Dirección del Trabajo. & Código del Trabajo Art. 25 bis; D. Aguayo & Crítica \\
\textbf{REQ-NEG-24} & Huella de Carbono GLEC & Computar y reportar de manera mensual las emisiones de gases de efecto invernadero ($\text{g CO}_2\text{e}/\text{t-km}$) auditables bajo norma GLEC e ISO 14083 para responder a las exigencias 2029 del cliente exportador (19\%). & Caso 10, Cap. 4.6; Entrevista A. Lecaros & Crítica \\
\textbf{REQ-NEG-25} & Trazabilidad Cliente 19\% & Proveer un portal web seguro para clientes que exponga en tiempo real la posición georreferenciada de la carga, temperatura del furgón y estado del despacho durante el tránsito del flete contratado. & Caso 10, Cap. 4.6; Entrevista A. Lecaros & Alta \\
\end{tabla}

\clearpage

% ============================================================================
%  ANEXO 2.B
% ============================================================================
\section*{Anexo 2.B: Inventario detallado de flota y caracterización de conductores}\label{sec:2-anexo-2b}
\addcontentsline{toc}{section}{Anexo 2.B: Inventario detallado de flota y caracterización de conductores}

Este anexo desglosa la infraestructura vehicular móvil y la fuerza laboral que compone la operación de Transportes Curimón S.A., diferenciando el régimen de propiedad, el nivel de equipamiento telemático basal y la estrategia de integración tecnológica requerida para cada segmento.

\subsection*{2.B.1 Inventario desagregado del parque vehicular (374 tractocamiones)}

El parque de tractocamiones se clasifica de acuerdo con su titularidad jurídica y su madurez instrumental en la \tabref{2-tractos-clasif}:

\begin{tabla}[ancho=completo]{Inventario clasificado de tractocamiones}{2-tractos-clasif}{L{38mm} C{16mm} C{20mm} C{20mm} L{36mm} X}
\textbf{Segmento de Flota} & \textbf{Cant.} & \textbf{Partic.} & \textbf{Antigüedad} & \textbf{Equipamiento Telemático Actual} & \textbf{Requerimiento de Integración} \\ \midrule
Flota Propia CAN bus Fábrica & 61 & 16,3\% & 3,2 años & Módulo telemático de fábrica con bus CAN inactivo. Nunca consultado. & Captura pasiva no intrusiva mediante acopladores inductivos sin corte de cableado. \\
Flota Propia sin Telemetría Fábrica & 87 & 23,3\% & 8,6 años & Sin telemetría de bus de datos. Dispositivos GPS básicos de primera generación. & Instalación de dispositivo telemático estándar con acoplamiento inductivo no invasivo. \\
Flota Terceros con GPS Previo & 192 & 51,3\% & Variable (4-12 a) & Dispositivos GPS de 3 proveedores comerciales dispares (Wialon, Wisetrack, Webfleet). & Ingesta estandarizada mediante interfaces API/Webhooks en modalidad de intercambio telemático. \\
Flota Terceros sin Dispositivo GPS & 34 & 9,1\% & Variable (>10 a) & Cero equipamiento tecnológico. Monitoreo puramente telefónico por voz. & Equipamiento telemático estándar e integración de posicionamiento georreferenciado. \\ \midrule
\textbf{Total Parque Tractocamiones} & \textbf{374} & \textbf{100,0\%} & \textbf{6,4 años (med)} & \textbf{Parque altamente heterogéneo.} & \textbf{Integración unificada bajo modelo agnóstico de ingesta telemática.} \\
\end{tabla}

\subsection*{2.B.2 Inventario de semirremolques y equipos de arrastre propios (210 unidades)}

Curimón es propietaria del 100\% de los 210 semirremolques utilizados en la operación, asegurando el acople físico de la carga independientemente de si el tracto motriz es propio o subcontratado (\tabref{2-semirremolques}):

\begin{tabla}[ancho=completo]{Inventario de semirremolques propios por tipología de carga}{2-semirremolques}{L{42mm} C{16mm} C{20mm} L{40mm} X}
\textbf{Tipología de Semirremolque} & \textbf{Cant.} & \textbf{\% Arrastre} & \textbf{Operación Principal y Clientes} & \textbf{Instrumental y Requerimiento Específico} \\ \midrule
Semirremolques Refrigerados (Reefers) & 44 & 21,0\% & Cadena de frío agroexportadora (diciembre-abril) e industria acuícola en Puerto Montt. & Sensorización con sondas Pt100 (-30 °C a +30 °C) y sensor de apertura de puertas. \\
Semirremolques Sustancias Peligrosas & 18 & 8,6\% & Transporte químico, combustibles y reactivos mineros (Antofagasta y Concepción). & Certificación bajo D.S. N.° 298/1994, rotulación NCh 2190, extintores y revisión especial. \\
Ramplas Planas y Furgones Secos & 98 & 46,7\% & Carga general paletizada, materiales industriales, retail y productos de consumo masivo. & Monitoreo de enganche, control de precintos y odometría de eje mediante sensor de rueda. \\
Tolvas Graneleras & 32 & 15,2\% & Transporte de graneles minerales, áridos y productos agrícolas a granel. & Control de accionamiento hidráulico de volteo y sensor de sobrepeso por eje. \\
Portacontenedores Portuarios & 18 & 8,5\% & Operación marítimo-portuaria en Valparaíso y San Antonio; contenedores de exportación. & Control de trabas de seguridad (\textit{twistlocks}) e integración documental aduanera. \\ \midrule
\textbf{Total Semirremolques Propios} & \textbf{210} & \textbf{100,0\%} & \textbf{Capacidad total de arrastre corporativo.} & \textbf{Flota 100\% de propiedad de Curimón S.A.} \\
\end{tabla}

\subsection*{2.B.3 Caracterización detallada de la dotación de conductores (454 operadores)}

La operación de la flota requiere una fuerza laboral de 454 choferes, estructurada en dos grupos con relaciones contractuales distintas (\tabref{2-conductores}):

\begin{tabla}[ancho=completo]{Caracterización sociolaboral y de control de conductores}{2-conductores}{L{40mm} X X}
\textbf{Atributo / Parámetro} & \textbf{Conductores Propios de Curimón} & \textbf{Conductores Subcontratados de Terceros} \\ \midrule
\textbf{Dotación Asignable} & \textbf{196 conductores (43,2\% de la fuerza laboral)} & \textbf{258 conductores (56,8\% de la fuerza laboral)} \\
\textbf{Relación Contractual} & Contrato de trabajo indefinido directo con Curimón S.A. & Dependientes laborales de los 148 transportistas terceros. Sin subordinación directa. \\
\textbf{Marco Jurídico Laboral} & Artículo 25 bis del Código del Trabajo (régimen especial). & Las obligaciones de Curimón y los transportistas dependen de la relación contractual y de los supuestos legales aplicables; verificar con asesoría laboral. \\
\textbf{Límites de Jornada Legal} & Máx. 5 h conducción continua; 2 h descanso; 8 h descanso diario; 180 h mensuales. & Mismos límites legales, pero con fiscalización histórica nula por parte de Curimón. \\
\textbf{Régimen de Turnos} & Controlado administrativamente por la empresa; turnos conocidos. & Desconocimiento total de turnos previos realizados para otros mandantes o clientes ajenos. \\
\textbf{Mecanismo Probatorio} & Registro digital inalterable de jornada con trazabilidad y custodia. & Atestación documental al aceptar despacho y registro de eventos operacionales del viaje. \\
\textbf{Resguardo de Privacidad} & Datos laborales sujetos a reglas de finalidad, acceso, seguridad y conservación. & Definir base jurídica y límites de tratamiento para cada etapa; considerar la Ley N.° 19.628 vigente a la fecha de la oferta y preparar la adaptación a la Ley N.° 21.719 (vigencia prevista: 1 de diciembre de 2026). \\
\textbf{Representante / Referente} & Yasna Colipán Marín (conductora de ruta norte, 7 años antigüedad). & Nolberto Sandoval Pinto (transportista con 2 tractos y chofer a cargo, 9 años en Curimón). \\
\end{tabla}

\clearpage

% ============================================================================
%  ANEXO 2.C
% ============================================================================
\section*{Anexo 2.C: Matriz exhaustiva de restricciones operacionales, legales y exclusiones contractuales}\label{sec:2-anexo-2c}
\addcontentsline{toc}{section}{Anexo 2.C: Matriz exhaustiva de restricciones operacionales, legales y exclusiones contractuales}

Este anexo recopila de forma sistemática las restricciones legales vigentes en Chile, las restricciones físicas de la operación en carretera y las exclusiones explícitas de la propuesta técnica.

\begin{tabla}[ancho=completo]{Matriz de restricciones legales, normativas y regulatorias}{2-rest-legales}{C{24mm} L{32mm} X X}
\textbf{Código} & \textbf{Norma / Estándar} & \textbf{Exigencia Normativa Concreta} & \textbf{Impacto en Diseño Operacional} \\ \midrule
\textbf{REST-LEG-01} & Código del Trabajo Art. 25 bis & A la fecha de la oferta: hasta 180 horas mensuales distribuidas en no menos de 21 días; descanso mínimo ininterrumpido de 8 horas en cada período de 24 horas; máximo de 5 horas de conducción continua seguido por al menos 2 horas de descanso. & Definir controles según el régimen aplicable al conductor y al momento de operación, con revisión laboral antes de implementarlos. \\
\textbf{REST-LEG-02} & Ley N.° 20.123 (Subcontratación) & La responsabilidad de la empresa principal se determina según los supuestos y condiciones establecidos en la legislación laboral aplicable. & Identificar con asesoría laboral los antecedentes y controles pertinentes para la relación entre Curimón y los transportistas externos. \\
\textbf{REST-LEG-03} & Ley N.° 21.377 y Ley de Tránsito & La regulación restringe la manipulación de dispositivos durante la conducción, con las condiciones y excepciones que establezcan la ley y su reglamentación vigente. & Evaluar controles de interfaz y voz para mitigar distracciones; validar excepciones y comportamiento con revisión normativa y pruebas de seguridad. \\
\textbf{REST-LEG-04} & Protección de datos personales & A la fecha de la oferta rige la Ley N.° 19.628; la Ley N.° 21.719 entra en vigor el 1 de diciembre de 2026. & Definir finalidad, base jurídica, acceso, conservación y medidas de seguridad conforme al régimen aplicable en cada etapa; validar con asesoría especializada. \\
\textbf{REST-LEG-05} & D.S. N.° 298/1994 (MTT - SUSPEL) & Reglamenta transporte de cargas peligrosas: curso de capacitación vigente, antigüedad máxima, revisión periódica, rotulado NCh 2190 y HDS. & Validación cruzada bloqueante en pre-despacho: verificación digital de la credencial del chofer y rótulos antes de destrabar orden de carga. \\
\textbf{REST-LEG-06} & D.S. N.° 43/2015 (MINSAL - Almacenamiento) & Regula el almacenamiento de sustancias peligrosas en instalaciones y actividades que estén dentro de su ámbito de aplicación. & Determinar aplicabilidad y controles para cada recinto y actividad con el responsable de prevención y la normativa vigente. \\
\textbf{REST-LEG-07} & D.S. N.° 158/1980 (MOP - Pesajes) & Límites máximos de peso bruto vehicular y peso por eje para preservar la infraestructura vial nacional. Control en plazas oficiales. & Monitoreo y control del manifiesto de estiba en origen para abatir las 142 detenciones registradas en 2025 (2.556 horas-camión inmovilizadas). \\
\textbf{REST-LEG-08} & Marco GLEC / ISO 14083:2023 & Cálculo estandarizado de emisiones de gases de efecto invernadero ($\text{g CO}_2\text{e}/\text{t-km}$) considerando factores Well-to-Wheel. & Reportabilidad auditada obligatoria para la renovación contractual de 2029 con el cliente exportador principal (19\% de ingresos corporativos). \\
\end{tabla}

\begin{tabla}[ancho=completo]{Matriz de restricciones físicas, ambientales y operacionales de la red}{2-rest-fisicas}{C{24mm} L{32mm} X X}
\textbf{Código} & \textbf{Restricción} & \textbf{Descripción del Límite Físico o Ambiental} & \textbf{Consecuencia en la Arquitectura} \\ \midrule
\textbf{REST-FIS-01} & Ventana de Intervención & La flota rueda 24/7/365. Los propios pasan por taller cada 6 días; el 22\% de los terceros pasa menos de 1 vez/mes por terminal. & Prohibición de detenciones masivas de flota. Despliegue secuencial de hardware en San Bernardo a razón de 25 camiones/mes. \\
\textbf{REST-FIS-02} & Sombra Celular Prolongada & Tramos de más de 80 kilómetros continuos sin cobertura de red móvil en el Desierto de Atacama (Ruta 5 Norte) y cordillera. & Arquitectura \textit{Offline-First}: persistencia en búfer local vehicular con retransmisión asíncrona determinista. \\
\textbf{REST-FIS-03} & Cierres Paso Los Libertadores & El caso informa cierres climáticos prolongados en alta montaña. & Mantener el mínimo de 72 horas de retención local exigido por RT-03.10; evaluar y justificar aparte cualquier capacidad superior propuesta para el dispositivo. \\
\textbf{REST-FIS-04} & Condiciones Térmicas y Vibratorias & Temperaturas extremas de cabina entre -20 °C y +70 °C y vibración severa continua en caminos pavimentados y no pavimentados. & Hardware certificado obligatoriamente bajo norma industrial SAE J1455 con sellado ambiental grado IP67. \\
\textbf{REST-FIS-05} & Infraestructura Servidores Matriz & Sala de servidores en San Bernardo de 26 m², split doméstico y UPS 20 min; incumple requerimientos RT-06.01 a RT-06.09. & Descarte de infraestructura central transaccional \textit{on-premise}; emplazamiento en nube pública resiliente (Azure Multi-AZ). \\
\textbf{REST-FIS-06} & Puntos de Clientes Ajenos (~1.400) & Recintos privados de terceros donde está estrictamente prohibido instalar infraestructura física, antenas o sensores de Curimón. & Trazabilidad basada exclusivamente en geocercas satelitales virtuales poligonales y telemetría propia del camión. \\
\end{tabla}

\begin{tabla}[ancho=completo]{Matriz de exclusiones contractuales explícitas}{2-exclusiones}{C{20mm} L{45mm} X}
\textbf{Código} & \textbf{Objeto Excluido} & \textbf{Justificación Técnica y Delimitación de Responsabilidad} \\ \midrule
\textbf{EXC-01} & Suministro de combustible diésel y lubricantes & La provisión física del carburante corresponde a los convenios de Curimón con distribuidoras mayoristas; audIT audita el consumo. \\
\textbf{EXC-02} & Mantenimiento mecánico y repuestos de taller & El recambio físico de piezas, neumáticos y reparaciones mecánicas es potestad exclusiva del personal de talleres de Curimón. \\
\textbf{EXC-03} & Modificación del código fuente del ERP contable & Curimón mantiene la titularidad de su software contable; audIT interoperará exclusivamente mediante servicios web y APIs. \\
\textbf{EXC-04} & Obras civiles mayores en terminales & Se excluyen ampliaciones edilicias en San Bernardo; la plataforma transaccional de misión crítica se alojará en la nube. \\
\textbf{EXC-05} & Honorarios y costos en Oferta Técnica & En cumplimiento del Artículo 50.2 de las Bases Administrativas, la propuesta técnica no contiene montos ni tarifas de audIT SpA. \\
\end{tabla}

\clearpage

% ============================================================================
%  ANEXO 2.D
% ============================================================================
\section*{Anexo 2.D: Fichas detalladas de caracterización de los trece (13) actores del ecosistema}\label{sec:2-anexo-2d}
\addcontentsline{toc}{section}{Anexo 2.D: Fichas detalladas de caracterización de los trece (13) actores del ecosistema}

A continuación se presentan las fichas completas de caracterización de los trece grupos de interés, integrando los tres actores omitidos en informes previos (Fondo de Inversión, Dirección del Trabajo y Aseguradora de Carga y Flota) y profundizando en las dependencias y riesgos de cada uno.

\subsection*{Ficha N.° 01: Fondo de Inversión Institucional (Accionista Minoritario 22\%)}
\begin{itemize}
  \item \textbf{Identificación:} Fondo de inversión institucional que adquirió el 22\% accionario en 2019. Posee asiento formal en Directorio.
  \item \textbf{Objetivos Estratégicos:} Maximización del retorno sobre el capital invertido (ROIC), resguardo del margen EBITDA, saneamiento de contratos deficitarios y mitigación de pasivos contingentes laborales.
  \item \textbf{Dolores Operacionales:} Margen comprimido al 9,0\%, 3 contratos principales bajo costo (-14\%) y contingencias laborales por Ley N.° 20.123.
  \item \textbf{Testimonio en Directorio:} \textit{«Exigimos auditoría de costos en tiempo real y gobernanza técnica certificada para proteger el valor de la compañía».}
  \item \textbf{Poder / Veto e Interés:} Poder Máximo en Directorio. Cuadrante 2: Mantener Satisfecho (Gobernanza y Retorno).
  \item \textbf{Mecanismo de Mitigación:} Reportabilidad mensual de indicadores consolidados de rentabilidad marginal por cliente y auditoría de riesgos bajo ISO 27001 e ISO 9001.
\end{itemize}

\subsection*{Ficha N.° 02: Dirección del Trabajo (DT --- Autoridad Laboral Fiscalizadora)}
\begin{itemize}
  \item \textbf{Identificación:} Órgano fiscalizador del Estado de Chile dependiente del Ministerio del Trabajo y Previsión Social.
  \item \textbf{Objetivos Estratégicos:} Fiscalización del cumplimiento estricto del régimen de jornada especial y descansos (Art. 25 bis y Ley N.° 20.123).
  \item \textbf{Dolores Operacionales:} No se registran descargas históricas de tacógrafos digitales en Curimón; además, existe información incompleta sobre jornadas previas de terceros y antecedentes de siniestros asociados a fatiga.
  \item \textbf{Testimonio en Fiscalización:} \textit{«Si no exhiben registros electrónicos automatizados e inalterables en la fiscalización, aplicaremos las multas máximas y clausura temporal».}
  \item \textbf{Poder / Veto e Interés:} Extremo (Potestad Pública Sancionatoria). Cuadrante 2: Mantener Satisfecho (Cumplimiento Legal Invariable).
  \item \textbf{Medida a evaluar:} Registro electrónico de jornada con controles de integridad y trazabilidad; confirmar periodicidad, formatos y exigencias de descarga de tacógrafos conforme a la regulación aplicable.
\end{itemize}

\subsection*{Ficha N.° 03: Compañías Aseguradoras de Carga y Flota}
\begin{itemize}
  \item \textbf{Identificación:} Entidades financieras aseguradoras que suscriben las pólizas de Responsabilidad Civil, carga perecible, SUSPEL y casco de tractocamiones.
  \item \textbf{Objetivos Estratégicos:} Determinación precisa del riesgo asegurable, exigencia de debida diligencia patronal y rechazo de coberturas ante negligencia inexcusable.
  \item \textbf{Dolores Operacionales:} 4 siniestros graves en 3 años, falta de telemetría continua de temperatura en reefers y ausencia de medios probatorios cinemáticos.
  \item \textbf{Testimonio de Ajustador:} \textit{«Si hay baches de temperatura en frío o el chofer manejó más de cinco horas, la cobertura caduca por negligencia patronal».}
  \item \textbf{Poder / Veto e Interés:} Alto (Denegación de Cobertura y Elevación de Primas). Cuadrante 2: Mantener Satisfecho (Mitigación Probatoria de Riesgo).
  \item \textbf{Mecanismo de Mitigación:} Sensorización Pt100 certificada en 44 ramplas de frío y registro circular local inalterable de telemetría cinemática a bordo.
\end{itemize}

\subsection*{Ficha N.° 04: Directorio y Familia Fundadora (78\% de la Propiedad)}
\begin{itemize}
  \item \textbf{Identificación:} Segunda generación familiar fundadora de Transportes Curimón S.A., controladora del 78\% del capital social.
  \item \textbf{Objetivos Estratégicos:} Continuidad histórica de la empresa, preservación del patrimonio familiar y defensa de la reputación de la marca Curimón.
  \item \textbf{Dolores Operacionales:} Pérdida de control operacional, amenaza por posible no renovación del contrato 19\% en 2029 y fricción con socios del fondo.
  \item \textbf{Testimonio Accionista:} \textit{«Queremos que la empresa sea viable para la tercera generación recuperando el control real en ruta».}
  \item \textbf{Poder / Veto e Interés:} Máximo Societario. Cuadrante 2: Mantener Satisfecho (Alineamiento Patrimonial y Reputacional).
  \item \textbf{Mecanismo de Mitigación:} Reportes trimestrales de avance en la erradicación de brechas normativas y cumplimiento del pliego licitado.
\end{itemize}

\subsection*{Ficha N.° 05: Gerencia General (María Teresa González)}
\begin{itemize}
  \item \textbf{Identificación:} Máxima autoridad ejecutiva de Curimón S.A., responsable ante el Directorio del desempeño financiero y operativo integral.
  \item \textbf{Objetivos Estratégicos:} Rentabilizar la operación, resolver la brecha tecnológica 2029 del cliente 19\% y asegurar un clima laboral estable.
  \item \textbf{Dolores Operacionales:} Asimetría de información entre operaciones y finanzas, falta de visibilidad en tiempo real de la flota y dependencia de despachos informales.
  \item \textbf{Poder / Veto e Interés:} Máximo Ejecutivo. Cuadrante 1: Gestionar Atentamente (Sponsor Principal del Proyecto).
  \item \textbf{Mecanismo de Mitigación:} Cuadro de mando integral en tiempo real con indicadores operacionales, financieros y de cumplimiento regulatorio.
\end{itemize}

\subsection*{Ficha N.° 06: Gerencia de Operaciones (Roberto Mansilla)}
\begin{itemize}
  \item \textbf{Identificación:} Responsable directo de la programación, despacho, seguimiento y cumplimiento de los itinerarios de los 374 camiones.
  \item \textbf{Objetivos Estratégicos:} Maximización de utilización de flota, reducción de los 10,66M de km en vacío y cumplimiento de ventanas horarias de clientes.
  \item \textbf{Dolores Operacionales:} Fricción continua con choferes por asignaciones manuales, falta de trazabilidad en los 226 camiones de terceros y desbalance de carga.
  \item \textbf{Poder / Veto e Interés:} Alto Ejecutivo Operacional. Cuadrante 1: Gestionar Atentamente (Usuario Líder de la Solución).
  \item \textbf{Mecanismo de Mitigación:} Módulo automatizado de despacho con optimización de retorno en vacío y visibilidad integral de flota en mapa único.
\end{itemize}

\subsection*{Ficha N.° 07: Jefatura de Prevención de Riesgos (Denise Aguayo)}
\begin{itemize}
  \item \textbf{Identificación:} Encargada de la seguridad laboral, fiscalización de jornadas, cumplimiento de D.S. 298 (SUSPEL) y vigencia documental de la flota.
  \item \textbf{Objetivos Estratégicos:} Cero fatalidades en ruta, cumplimiento irrestricto del Art. 25 bis y erradicación de multas laborales y de tránsito.
  \item \textbf{Dolores Operacionales:} Gestión manual inmanejable de ~6.000 vigencias en planillas Excel y falta de medios para fiscalizar a los 258 choferes terceros.
  \item \textbf{Poder / Veto e Interés:} Alto (Potestad Interna de Bloqueo de Despacho). Cuadrante 1: Gestionar Atentamente.
  \item \textbf{Mecanismo de Mitigación:} Validación automatizada bloqueante de vigencias en pre-despacho y alertas pasivas de fatiga en ruta.
\end{itemize}

\subsection*{Ficha N.° 08: Gerencia de Finanzas y Control de Gestión (Gabriel Ossandón)}
\begin{itemize}
  \item \textbf{Identificación:} Responsable de la facturación, cobranza, costeo marginal y liquidación de fletes a los 148 transportistas terceros.
  \item \textbf{Objetivos Estratégicos:} Cobro efectivo de sobreestadías (recuperar el 71\% objetado), costeo real por viaje y aceleración del ciclo de facturación.
  \item \textbf{Dolores Operacionales:} Pérdida anual de ingresos por falta de sustento probatorio en andenes y 9 días de demora en pre-liquidaciones con 11\% de error.
  \item \textbf{Poder / Veto e Interés:} Alto Financiero. Cuadrante 1: Gestionar Atentamente.
  \item \textbf{Mecanismo de Mitigación:} Registro cronológico digital certificado de estadías en clientes y módulo automatizado de pre-liquidación a terceros.
\end{itemize}

\subsection*{Ficha N.° 09: Jefatura de Mantenimiento y Talleres (Hernán Trincado)}
\begin{itemize}
  \item \textbf{Identificación:} Responsable del mantenimiento preventivo y correctivo de los 148 tractocamiones y 210 semirremolques propios en taller San Bernardo.
  \item \textbf{Objetivos Estratégicos:} Maximizar disponibilidad mecánica de flota propia, prevenir panas en ruta y auditar consumos anómalos de diésel.
  \item \textbf{Dolores Operacionales:} Planificación basada en pautas fijas de 6 días en lugar de desgaste real y desaprovechamiento del bus CAN en 61 unidades.
  \item \textbf{Poder / Veto e Interés:} Medio Operacional. Cuadrante 3: Mantener Informado.
  \item \textbf{Mecanismo de Mitigación:} Telemetría CAN bus pasiva continua para generación automática de órdenes de trabajo por odometría y horas de motor.
\end{itemize}

\subsection*{Ficha N.° 10: Jefatura de Tecnologías de Información (Patricia Kast)}
\begin{itemize}
  \item \textbf{Identificación:} Responsable de la infraestructura informática, bases de datos locales, ERP contable y enlaces de telecomunicaciones de Curimón.
  \item \textbf{Objetivos Estratégicos:} Modernización hacia arquitectura cloud escalable, alta disponibilidad del sistema y seguridad de la información.
  \item \textbf{Dolores Operacionales:} Sala de servidores precaria en San Bernardo (26 m², split doméstico, UPS 20 min) y multiplicidad de sistemas no integrados.
  \item \textbf{Poder / Veto e Interés:} Medio Técnico. Cuadrante 1: Gestionar Atentamente (Contraparte Técnica de Implantación).
  \item \textbf{Mecanismo de Mitigación:} Despliegue en Microsoft Azure Multi-AZ con API Gateway, base de datos relacional gestionada y soporte nivel 3 continuo.
\end{itemize}

\subsection*{Ficha N.° 11: Despachadores y Operadores de Torre de Control (Esteban Valdebenito)}
\begin{itemize}
  \item \textbf{Identificación:} Equipo de 22 despachadores distribuidos en los 5 terminales que coordinan la salida física de camiones y monitoreo radial.
  \item \textbf{Objetivos Estratégicos:} Liberación ágil de cargas, resolución rápida de contingencias en ruta y atención oportuna de llamadas de clientes.
  \item \textbf{Dolores Operacionales:} Sobrecarga administrativa, dispersión de información en 3 pantallas de GPS y llamadas telefónicas continuas a choferes.
  \item \textbf{Poder / Veto e Interés:} Medio Operacional Directo. Cuadrante 3: Mantener Informado / Capacitar Intensivamente.
  \item \textbf{Mecanismo de Mitigación:} Consola unificada de despacho con semáforo de aptitud vehicular y mapas de seguimiento en tiempo real.
\end{itemize}

\subsection*{Ficha N.° 12: Conductores Propios de Curimón (Yasna Colipán)}
\begin{itemize}
  \item \textbf{Identificación:} Dotación de 196 choferes contratados directamente por Curimón, con experiencia en rutas troncales norte y sur.
  \item \textbf{Objetivos Estratégicos:} Conducción segura, respeto irrestricto de descansos legales, condiciones dignas de viaje y retribución justa sin descuentos injustificados.
  \item \textbf{Dolores Operacionales:} Fatiga acumulada por llamadas en ruta, presiones comerciales para apurar tránsitos y temor a invasión de privacidad.
  \item \textbf{Poder / Veto e Interés:} Alto Sindical / Operacional. Cuadrante 1: Gestionar Atentamente (Adopción de Usuarios Finales).
  \item \textbf{Mecanismo de Mitigación:} Enclavamiento cinético de pantallas (Ley No Chat), alertas de descanso pasivas con anticipación y cero cámaras invasivas en cabina.
\end{itemize}

\subsection*{Ficha N.° 13: Transportistas Subcontratados y Terceros (Nolberto Sandoval)}
\begin{itemize}
  \item \textbf{Identificación:} 148 pequeños y medianos empresarios dueños de los 226 camiones subcontratados y empleadores de 258 choferes externos.
  \item \textbf{Objetivos Estratégicos:} Cobro transparente y puntual de fletes, autonomía en sus activos vehiculares y acceso a tarifas justas y estables.
  \item \textbf{Dolores Operacionales:} Retraso de 9 días en liquidaciones con errores frecuentes, desconfianza ante exigencias de instalar hardware y temor a fiscalización fuera de viaje.
  \item \textbf{Poder / Veto e Interés:} Alto Capacidad de Bloqueo Operacional (60,4\% de la flota). Cuadrante 1: Gestionar Atentamente.
  \item \textbf{Mecanismo de Mitigación:} Pre-liquidaciones automáticas transparentes, homologación de GPS existentes vía API y desconexión obligatoria de telemetría fuera de orden de flete.
\end{itemize}
